\chapter{Lanes, short lanes and movements}
\section{Regular lanes}
\begin{lstlisting}[language=yaml]
lanes:
  - id: A1_1
    intersection_id: I1
    physical_length_m: 220
    saturation_flow:
      method: direct
      value_pcu_per_hour: 1900
\end{lstlisting}
If storage is derived from physical length, the basic relationship is
\begin{equation}
K_{\mathrm{lane}}=\frac{L_{\mathrm{lane}}}{l_{\mathrm{PCU}}}.
\end{equation}
With a 7.35 m PCU, a 220 m lane stores approximately 29.9 PCU.

\section{Short lanes}
A short lane references an existing parent lane. Multiple short lanes may share the same parent.
\begin{lstlisting}[language=yaml]
  - id: A2_3L
    intersection_id: I1
    parent_lane_id: A2_3
    physical_length_m: 80
    saturation_flow:
      method: rr67_unopposed
      lane_width_m: 3.2
\end{lstlisting}

\section{Movements are lane-to-lane paths}
A movement should be read as
\begin{equation}
\textit{movement}=\textit{from lane}\rightarrow\textit{to lane}.
\end{equation}
\begin{lstlisting}[language=yaml]
movements:
  - id: T_WE
    intersection_id: I1
    from_lane_id: A1_1
    to_lane_id: A3_4
    movement_type: through
    signal_group_id: SG_A1
    free_flow_time_s: 8
\end{lstlisting}
\code{to_lane_id} is optional for a movement that exits the model, but when supplied it should reference a real lane.

\section{Several movements on one lane}
If \code{T_WE} is 500 PCU/h and \code{L_WN} is 100 PCU/h and both use \code{A1_1}, the physical lane demand is 600 PCU/h. PySig retains movement composition but calculates one physical queue, delay process and storage constraint for \code{A1_1}.
